\section{Methodology}
This systematic review was conducted in accordance with the Preferred Reporting Items for Systematic Reviews and Meta-Analyses (PRISMA) 2020 statement guidelines \citep{PRISMA2020}. The review protocol was registered prospectively in the PROSPERO database (Registration ID: CRD420261093).

\subsection{Information Sources and Search Strategy}
A comprehensive electronic search was conducted across three databases: PubMed, Embase, and the Cochrane Central Register of Controlled Trials (CENTRAL). The search window spanned from database inception to July 1, 2026. A combination of Medical Subject Headings (MeSH) and free-text terms was utilized, including terms such as "deep brain stimulation", "closed-loop", "adaptive DBS", "treatment-resistant depression", and "neuromodulation". In addition, reference lists of relevant review articles and included studies were hand-searched to identify missed citations.

\subsection{Eligibility Criteria}
Studies were selected based on the following PICOS (Population, Intervention, Comparison, Outcome, Study design) criteria:
\begin{itemize}
  \item \textbf{Population:} Adult patients (aged $\ge 18$ years) diagnosed with unipolar or bipolar major depressive disorder who met the criteria for treatment-resistant depression (failure of at least two antidepressant trials and one electroconvulsive therapy trial).
  \item \textbf{Intervention:} Active therapy using an implanted closed-loop (adaptive or responsive) deep brain stimulation system.
  \item \textbf{Comparison:} Active vs. sham stimulation phases, pre-to-post implant baseline comparisons, or open-loop controls.
  \item \textbf{Outcomes:} Efficacy measured by change in validated depression scales (e.g., Montgomery-Asberg Depression Rating Scale (MADRS) or Hamilton Depression Rating Scale (HAMD)). Safety measured by the reporting of treatment-emergent adverse events (TEAEs).
  \item \textbf{Study Design:} Randomized controlled trials (RCTs), open-label clinical trials, prospective cohort studies, and case series with $N \ge 3$ participants.
\end{itemize}

\subsection{Study Selection and Data Extraction}
Two review authors independently screened titles and abstracts against the eligibility criteria. Full-text articles of potentially relevant citations were retrieved and assessed for eligibility. Any discrepancies were resolved via consensus or by consultation with a third reviewer. Data were extracted using a standardized form, including study demographics, device manufacturer (e.g., Synapse, Meridian), brain target, neural biomarkers, follow-up duration, response rate (defined as $\ge 50$\% reduction in depression score), and remission rate.

\subsection{Quality and Risk of Bias Assessment}
Risk of bias was evaluated using the Cochrane Risk of Bias tool for randomized trials (RoB 2) and the Risk of Bias in Non-randomized Studies of Interventions (ROBINS-I) tool. The quality of evidence was appraised across five domains: bias arising from the randomization process, bias due to deviations from intended interventions, bias due to missing outcome data, bias in measurement of the outcome, and bias in selection of the reported result.
